\subsection{Experimental algorithmics}
The case for measuring algorithms rather than only analysing them was made
systematically by Johnson~\cite{johnson2002theoretician}, whose guide for
theoreticians lists the pitfalls this paper tries to avoid (uninteresting
instances, unfair baselines, unreported machines), and by
McGeoch~\cite{mcgeoch2012guide}, whose book is the standard reference on
experimental design for algorithms. Sanders~\cite{sanders2009algorithm} frames
the activity as \emph{algorithm engineering}: a cycle of design, analysis,
implementation and experiment in which the constant factors are first-class
objects. The measurement side has its own literature of cautionary results.
Mytkowicz et al.~\cite{mytkowicz2009producing} showed that link order and
environment-variable size alone can change measured runtimes by tens of
percent, more than many published effects; Curtsinger and
Berger~\cite{curtsinger2013stabilizer} proposed randomising memory layout so
that such effects average out; and Hoefler and Belli~\cite{hoefler2015scientific}
surveyed reporting practice and argued for repetition, interval estimates and
full disclosure of the machine. LaMarca and Ladner~\cite{lamarca1999influence}
is the direct ancestor of this study for sorting: by instrumenting cache
misses they showed that the instruction count of a sorting algorithm is a poor
predictor of its time once the input leaves cache, and that memory-aware
variants of heapsort and mergesort outrun the textbook versions.

\subsection{The four pairs}
\paragraph{Sorting.}
Introsort~\cite{musser1997introspective}---quicksort with a heapsort fallback
and insertion sort for small ranges---is what \texttt{std::sort} runs in every
major C++ standard library, which makes it the baseline a practitioner actually
faces. Rahman and Raman~\cite{rahman2001adapting} adapted LSD radix sort to the
cache and TLB and showed it beating comparison sorts on large arrays of
integers; Satish et al.~\cite{satish2010fast} compared radix and merge sort
across CPUs and GPUs and found radix sort faster for 32-bit keys but merge sort
preferable for wide keys and on bandwidth-limited hardware; Polychroniou and
Ross~\cite{polychroniou2014comprehensive} studied the partitioning step that
both families share at billion-key scale. None of these report the crossover
size against \texttt{std::sort} as an estimated quantity, and all use SIMD
and multithreading that move it.

\paragraph{Matrix multiplication.}
Strassen's 1969 algorithm~\cite{strassen1969gaussian} is the archetype of an
asymptotic improvement that is practical only after engineering.
Huss-Lederman et al.~\cite{husslederman1996implementation} gave the first
careful implementation study, establishing that the recursion must stop at a
tuned cutoff and that odd dimensions need peeling or padding; D'Alberto and
Nicolau~\cite{dalberto2007adaptive} tuned that cutoff adaptively to the memory
hierarchy; Benson and Ballard~\cite{benson2015framework} showed that the best
fast algorithm depends on matrix shape; and Huang et
al.~\cite{huang2016strassen}, by fusing Strassen's additions into a BLIS
kernel, brought the \emph{switch point}---the dimension above which one level of
Strassen pays off against a vendor-grade DGEMM---down from the roughly 2\,000
of earlier implementations to as little as 500 for their best variant, on a
single Ivy Bridge Xeon core. Those crossovers are against the strongest
possible baseline. Ours is against the baseline a
programmer writes without a BLAS, and is correspondingly much earlier.

\paragraph{Integer multiplication.}
Karatsuba and Ofman's algorithm~\cite{karatsuba1962multiplication} is the
first sub-quadratic multiplication and the first step of the tower that ends,
for now, at Harvey and van der Hoeven's $\bigO(n\log n)$
algorithm~\cite{harvey2021integer}, which is correct and
unimplementable. Bernstein~\cite{bernstein2001multidigit} gives the modern
unified view. The GNU Multiple Precision library~\cite{gmp2023manual} switches
from schoolbook to Karatsuba at a machine-tuned threshold
(\texttt{MUL\_TOOM22\_THRESHOLD}), which its manual notes can be as low as ten
64-bit limbs; the parameter file GMP ships for Alder Lake sets it to 13
limbs~\cite{gmp2022alderlake}. This gives an external reference point for our
measured crossover.

\paragraph{Shortest paths.}
Fredman and Tarjan~\cite{fredman1987fibonacci} introduced the Fibonacci heap
precisely to improve Dijkstra's algorithm from $\bigO(m\log n)$ to
$\bigO(m+n\log n)$, and the data structure has been a textbook fixture ever
since; Brodal, Lagogiannis and Tarjan~\cite{brodal2012strict} later made the
bounds worst-case. The empirical record is uniformly unkind to it. Cherkassky,
Goldberg and Radzik~\cite{cherkassky1996shortest} found simple heaps and bucket
structures faster on every graph family they tried; Chen et
al.~\cite{chen2007priority} dropped the Fibonacci heap from their comparison
after preliminary runs and found that Dijkstra variants \emph{without}
decrease-key were at least $1.4\times$ faster than those with it on random
graphs of average degree 8 (and faster on every road network they tried); and
Larkin, Sen and Tarjan~\cite{larkin2014backtobasics}, on the full USA road
network, measured the Fibonacci heap at $3.15\times$ the time of an implicit
4-ary heap, with one decrease-key for every fourteen inserts, and found the L1
cache-miss rate to be the counter most strongly correlated with wall-clock
time. Our results
agree with all of these and add the one regime, complete graphs, in which the
asymptotic argument finally holds.

\subsection{Hidden constants and galactic algorithms}
Lipton~\cite{lipton2010galactic} coined \emph{galactic algorithm} for
one whose asymptotic advantage appears only on inputs larger than any that
will be run, with fast matrix multiplication as the canonical example: the
exponent has been pushed to $\omega<2.372$~\cite{alman2024refined,
alman2024asymmetry}, yet nothing past Strassen is used in practice. This paper
sits at the opposite end of that spectrum. Its four pairs are algorithms that
are actually used, and the question is not whether their crossovers exist but
where, with what precision, and what moves them. The hardware background that
answers ``why''---cache levels, memory-level parallelism, branch
prediction---is laid out by Drepper~\cite{drepper2007memory} and Fog~\cite{fog2026optimizing}; the
Linux \texttt{perf\_event\_open} interface~\cite{linux2026perfeventopen} is
how the harness reads the counters where it can.
